\begin{figure}[ht]
\centering
\resizebox{\textwidth}{!}{%
\begin{tikzpicture}[
    font=\small,
    caja/.style={
        draw,
        rectangle,
        rounded corners=4pt,
        align=center,
        minimum height=1cm,
        text width=3.4cm,
        fill=white,
        inner sep=4pt
    },
    cajaGrande/.style={
        draw,
        rectangle,
        rounded corners=4pt,
        align=center,
        minimum height=1.2cm,
        text width=6cm,
        fill=white,
        inner sep=4pt
    },
    cajaCausa/.style={
        draw,
        rectangle,
        rounded corners=4pt,
        align=center,
        minimum height=1cm,
        text width=5cm,
        fill=white,
        inner sep=4pt
    },
    linea/.style={-{Stealth[length=4mm]}, thick}
]

% Efecto superior
\node[cajaGrande] (impacto) at (0,0)
    {Impacto educativo reducido \\del programa};

% Efectos intermedios
\node[caja] (exposicion) at (-9.0,-3.5)
    {Baja exposición\\ estudiantil};
\node[caja] (docentes) at (-4.5,-3.5)
    {Docentes con acceso limitado \\a programas};
\node[caja] (interes) at (0,-3.5)
    {Interés juvenil sin\\ espacios formales};
\node[caja] (alcance) at (4.5,-3.5)
    {Alcance \\organizativo \\restringido};
\node[caja] (continuidad) at (9.0,-3.5)
    {Continuidad reducida\\ del programa};

% Problema central
\node[cajaGrande] (problema) at (0,-7.5)
    {Limitaciones en la \\implementación del Desafío \\Bebras en Bolivia para \\estudiantes escolares};

% Causas intermedias
\node[cajaCausa] (infraestructura) at (-4.5,-11.5)
    {Infraestructura tecnológica\\ no especializada};
\node[cajaCausa] (gestion) at (4.5,-11.5)
    {Gestión fragmentada\\ del concurso};

% Causas raíz
\node[cajaCausa] (presencia) at (-4.5,-14.5)
    {Pensamiento computacional\\ con presencia marginal\\ en el aula};
\node[cajaCausa] (iniciativas) at (4.5,-14.5)
    {Iniciativas previas sin\\ continuidad};

% Flechas efectos intermedios -> efecto superior
\draw[linea] (exposicion.north) -- (impacto.south west);
\draw[linea] (docentes.north) -- (impacto.south west);
\draw[linea] (interes.north) -- (impacto.south);
\draw[linea] (alcance.north) -- (impacto.south east);
\draw[linea] (continuidad.north) -- (impacto.south east);

% Flechas problema central -> efectos intermedios (CORREGIDO)
\draw[linea] (problema.north west) -- (exposicion.south);
\draw[linea] (problema.north west) -- (docentes.south);
\draw[linea] (problema.north) -- (interes.south);
\draw[linea] (problema.north east) -- (alcance.south);
\draw[linea] (problema.north east) -- (continuidad.south);

% Flechas causas intermedias -> problema central
\draw[linea] (infraestructura.north) -- (problema.south west);
\draw[linea] (gestion.north) -- (problema.south east);

% Flechas causas raíz -> causas intermedias
\draw[linea] (presencia.north) -- (infraestructura.south);
\draw[linea] (iniciativas.north) -- (gestion.south);

\end{tikzpicture}%
}
\caption{Árbol del problema para la implementación del Desafío Bebras Bolivia}
\label{fig:arbol-problemas}
\end{figure}